\subsection{绝对正则代数或绝对正规代数}

定义 1.--- 设 \(k\) 为域，\(A\) 为 \(k\)-代数。称 \(A\) 是绝对正则的（相应地绝对正规的）\(^{1}\)，如果环 \(\mathrm{A}_{(k')}\) 对每个有限次纯不可分扩张 \(k'\)（在 \(k\) 上）都是正则的（相应地正规的）。

每个绝对正则（相应地绝对正规）的 \(k\)-代数都是正则的（相应地正规的），在定义 1 中取 \(k' = k\) 即可看出。回忆：\(k\)-代数 A 若使环 \(\mathrm{A}_{(k')}\) 对每个有限次纯不可分扩张 \(k'\)（在 \(k\) 上）都为约化，则它就是 \(k\)-可分代数 (A, V, 第 117 页, 定理 2)。

例.--- 1) 若 k 是完全域，则每个正则（相应地正规）的 k-代数都是绝对正则的（相应地绝对正规的）。

\begin{enumerate}
\def\labelenumi{\arabic{enumi})}
\setcounter{enumi}{1}
\tightlist
\item
  设 A 为 Artin \(k\)-代数。若 A 正规，则它约化，从而同构于 \(k\) 的一些扩张的有限积 (A, VIII, § 8, 第 1 节, 命题 2)。下列条件等价：
\end{enumerate}

A 是可分的；

A 是绝对正则的；

A 是绝对正规的。

事实上，若 A 绝对正规，则它可分；若 A 可分，则对每个有限次纯不可分扩张 \(k'\)（在 \(k\) 上），环 \(\mathrm{A}_{(k')}\) 都是约化的且 Artin 的，从而同构于一些域的积，因而是正则的，于是 \(A\) 是绝对正则的。

此外，若 k-代数 A 是有限次的，则上述条件等价于说 A 是平展的 (A, V, 第 34 页, 定理 4)。

\begin{enumerate}
\def\labelenumi{\arabic{enumi})}
\setcounter{enumi}{2}
\tightlist
\item
  设 A 为正则局部 \(k\)-代数。若扩张 \(\kappa_{\mathrm{A}}\)（在 \(k\) 上）可分，则代数 A 绝对正则。事实上，设 \(k'\) 为 \(k\) 的一个有限扩张。A-模 \(\mathrm{A}_{(k')}\) 是自由的。由于它是有限型的，\(\mathrm{A}_{(k')}\) 的每个极大理想都位于 \(\mathfrak{m}_{\mathrm{A}}\) 之上 (V, § 2, 第 1 节, 命题 1)。环 \(\kappa_{\mathrm{A}} \otimes_{\mathrm{A}} \mathrm{A}_{(k')}\) 同构于 \(\kappa_{\mathrm{A}} \otimes_{k} k'\)，是正则的 (第 3 节, 引理 2)。由 § 4 第 5 节命题 9 的推论，环 \(\mathrm{A}_{(k')}\) 正则，这就证明了 \(k\)-代数 A 绝对正则。
\end{enumerate}

命题 6. 设 k 为域，A 为 Noether k-代数。

\begin{enumerate}
\def\labelenumi{\alph{enumi})}
\item
  若 A 绝对正则（相应地绝对正规、可分），则对 A 的每个乘性子集 S，\(S^{-1}A\) 亦然。
\item
  若对 A 的每个极大理想 m，\(A_{m}\) 都绝对正则（相应地绝对正规、可分），则 A 绝对正则（相应地绝对正规、可分）。
\item
  这由下述事实得出：环 \((\mathrm{S}^{-1}\mathrm{A})_{(k^{\prime})}\) 同构于环 \(\mathrm{A}_{(k^{\prime})}\) 的一个分式环，这对 k 的每个扩张 \(k^{\prime}\) 都成立。
\item
  设对 A 的每个极大理想 m，\(A_{m}\) 都绝对正则（相应地绝对正规、可分）。设 \(k'\) 为 k 的一个有限次纯不可分扩张，设 \(m'\) 为 \(\mathrm{A}_{(k')}\) 的一个极大理想。尚需验证局部环 \((\mathrm{A}_{(k')})_{\mathfrak{m}'}\) 是正则的（相应地正规的、约化的）。
\end{enumerate}

典范同态 \(A \to A_{(k')}\) 使 \(A_{(k')}\) 成为有限 A-代数。理想 \(m'\) 位于 A 的某个极大理想 m 之上 (V, § 2, 第 1 节, 命题 1)，而 \((\mathrm{A}_{(k')})_{\mathfrak{m}'}\) 同构于正则（相应地正规、约化）环 \((\mathrm{A}_{\mathfrak{m}})_{(k')}\) 的一个分式环，因而是正则的（相应地正规的、约化的）。

引理 3.--- 设 k 为域，K 为 k 的一个有限生成扩张。存在 K 的一个有限次扩张 L，以及 L 在 k 上的一个中间扩张 k′，它对 k 是纯不可分的且有限次，使得 L 在 k′ 上的扩张可分。

取定 K 与一个完全闭包 \(\hat{k}\)（在 \(k\) 上）的复合扩张 E，以及 E 的元素组 \((t_1, \ldots, t_n)\)，使得 E 是 \(\hat{k}(t_1, \ldots, t_n)\) 的可分代数扩张。设 I 为 K 在 \(k\) 上的一个有限生成系。域 \(\hat{k}(t_1, \ldots, t_n)\)（相应地 \(\hat{k}\mathrm{K}\)）是子域 \(k'(t_1, \ldots, t_n)\)（相应地 \(k'\mathrm{K}\)）的并，其中 \(k'\) 取遍 \(\hat{k}\) 的对 \(k\) 有限的中间扩张；因此可以找到这样一个中间扩张 \(k'\)，使得元素 \(t_1, \ldots, t_n\)（它们属于 \(\mathrm{E} = \hat{k}\mathrm{K}\)）含于 \(k'\mathrm{K}\)，且 I 的每个元素在 \(k'(t_1, \ldots, t_n)\) 上都是代数的且可分的。于是 \(\mathrm{L} = k'\mathrm{K}\) 是 \(k'\) 的可分扩张且对 K 有限次，而 \(k'\) 是 \(k\) 的有限次纯不可分扩张。

命题 7. 设 \(k\) 为域，A 与 B 为 \(k\)-代数，其中一个是本质有限型的。设 A 绝对正则（相应地绝对正规）。若环 B 正则（相应地正规），则环 A \(\otimes_{k}\) B 亦然。

设 K 为 k 的有限生成扩张；我们来证明环 \(\mathrm{A}_{(\mathrm{K})}\) 正则（相应地正规）。事实上，考察具有引理 3 中性质的扩张 L 与 \(k'\)。于是环 \(\mathrm{A}_{(k')}\) 依定义正则（相应地正规），而环 \(\mathrm{A}_{(\mathrm{L})}\)——它与 \(\mathrm{A}_{(k')} \otimes_{k'} \mathrm{L}\) 恒同——由第 3 节命题 5, b) 是正则的（相应地正规的）。于是由第 3 节命题 5, a)，\(\mathrm{A}_{(\mathrm{K})}\) 正则（相应地正规）。

设环 B 正则（相应地正规），我们来证明命题。典范同态 \(B \to A \otimes_{k} B\) 使 \(A \otimes_{k} B\) 成为自由 B-模。对 B 的每个素理想 p，环 \((A \otimes_{k} B) \otimes_{B} \kappa(\mathfrak{p})\) 由 § 4 第 5 节命题 9 的推论（相应地 § 1 第 10 节定理 4 的推论 3）与 \(A \otimes_{k} \kappa(\mathfrak{p})\) 恒同；我们只需证明对 B 的每个素理想 p，\(A \otimes_{k} \kappa(\mathfrak{p})\) 都正则（相应地正规）。

若 \(k\)-代数 B 本质有限型，则 \(\kappa(\mathfrak{p})\) 在 \(k\) 上的扩张是有限型的，且由上文所见，环 A \(\otimes_k \kappa(\mathfrak{p})\) 正则（相应地正规）。今设 \(k\)-代数 A 本质有限型；环 A \(\otimes_k \kappa(\mathfrak{p})\) 是 Noether 的，并且是 Noether 子环 \(A \otimes_k K\) 的递增有向族之并，其中 \(K\) 取遍 \(\kappa(\mathfrak{p})\) 的有限生成中间扩张。后述这些环都正则（相应地正规），于是可应用第 3 节的引理 1。

推论 1. 设 k 为域。两个绝对正则（相应地绝对正规）的 k-代数（其中一个是本质有限型的）的张量积是绝对正则（相应地绝对正规）的 k-代数。

设 A 与 B 为满足推论假设的两个 k-代数。设 \(k'\) 为 k 的一个有限次纯不可分扩张。环 \(\mathrm{B}_{(k')}\) 正则（相应地正规）；由命题 7，环 \(A \otimes_{k} B_{(k')}\) 正则（相应地正规），而与之同构的环 \((\mathrm{A} \otimes_{k} \mathrm{B}) \otimes_{k} k'\) 亦然。

推论 2.--- 设 k 为域，A 为绝对正则（相应地绝对正规）的 k-代数，K 为 k 的扩张。设 A 本质有限型，或 k 的扩张 K 有限型。

\begin{enumerate}
\def\labelenumi{\alph{enumi})}
\item
  环 \(\mathrm{A}_{(K)}\) 正则（相应地正规）。
\item
  若 k 的扩张 K 可分，则 k-代数 \(\mathrm{A}_{(\mathrm{K})}\) 绝对正则（相应地绝对正规）。
\end{enumerate}

断言 (a) 由命题 7 得出；断言 (b) 由推论 1 与例 2 得出。

推论 3. 设 \(k\) 为域，A 为 \(k\)-代数，K 为 \(k\) 的扩张。设 \(k\)-代数 A 本质有限型，或 \(k\) 的扩张 K 有限型。为使 \(k\)-代数 A 绝对正则（相应地绝对正规），必须且只需 K-代数 \(\mathrm{A}_{(\mathrm{K})}\) 绝对正则（相应地绝对正规）。

设 A 绝对正则（相应地绝对正规），并设 \(K'\) 为 K 的一个有限次纯不可分扩张。环 \(K'\otimes_{K}A_{(K)}\) 同构于 \(K'\otimes_{k}A\)，由推论 2 是正则的（相应地正规的）。

反之，设 K-代数 \(\mathrm{A}_{(\mathrm{K})}\) 绝对正则（相应地绝对正规），并设 \(k'\) 为 \(k\) 的一个有限次纯不可分扩张。设 L 为 \(k'\) 与 K 的一个复合扩张；于是环 \(\mathrm{A}_{(\mathrm{L})}\) 与 \(\mathrm{L} \otimes_{\mathrm{K}} \mathrm{A}_{(\mathrm{K})}\) 恒同，因而是正则的（相应地正规的）；从而由第 3 节命题 5, a)，环 \(\mathrm{A}_{(k')}\) 正则（相应地正规）。

推论 4.--- 设 k 为域，A 为本质有限型的 k-代数，K 为 k 的一个作为完全域的扩张。为使 A 绝对正则（相应地绝对正规），必须且只需环 \(\mathrm{A}_{(\mathrm{K})}\) 正则（相应地正规）。

这由推论 3 与例 1 得出。

